\section{Exams}

\subsection{2026}

\subsubsection{July 17}

\exercise Select all (and only those) statements correctly defining classification, regression, representation learning, and clustering.
\begin{enumerate}
    \item A model trained only to reconstruct its input is necessarily a reinforcement learning agent.
    \item A machine-health estimator may combine a discrete status label and a continuous risk score.
    \item Regression predicts continuous quantities such as price or temperature.
    \item Clustering requires the desired class label for every training example.
    \item A task is unsupervised only because its input is high-dimensional.
    \item Classification predicts one or more discrete labels.
    \item Predicting a binary sentiment label from examples is supervised learning.
    \item Learning an embedding can be useful even if no semantic class label is available.
\end{enumerate}
\solution Before looking at the answers, reduce everything to \textbf{4 definitions}:
\begin{itemize}
    \item \definition{Classification}
        \begin{center}
            \emph{Are we predicting a \textbf{discrete label/category}?}
        \end{center}
        We have an input $x$, and we want to predict a \textbf{discrete label}:
        \begin{equation*}
            f(x) = \hat{y}, \quad \hat{y} \in \left\{C_{1}, C_{2}, \ldots, C_{K}\right\}
        \end{equation*}
        For example, given an image, we want to predict whether it contains a cat, a dog, or a bird. The word \textbf{discrete} means that the output is a finite set of possible values, rather than a continuous range.

        \highspace
        Also, classification can be binary (two classes, e.g., spam vs. not spam) or multi-class (more than two classes, e.g., cat, dog, bird). In some cases, we may also have \textbf{multi-label classification}, where \hl{each input can belong to multiple classes simultaneously.}

    \item \definition{Regression}
        \begin{center}
            \emph{Are we predicting a \textbf{continuous numerical value}?}
        \end{center}
        Again, we have an input $x$, but now the target is a \textbf{continuous quantity}:
        \begin{equation*}
            f(x) = \hat{y}, \quad \hat{y} \in \mathbb{R}
        \end{equation*}
        For example, given the features of a house (size, number of bedrooms, location), we want to predict its price. The crucial \textbf{distinction} from classification is therefore \textbf{what the prediction represents}: a discrete category vs. a continuous value.

    \item \definition{Clustering}
        \begin{center}
            \emph{Are we finding \textbf{groups in unlabeled data}?}
        \end{center}
        This is different from classification and regression. We are \textbf{not given the correct group for every example}. Instead, we have data:
        \begin{equation*}
            D = \left\{x_1, x_2, \ldots, x_N\right\}
        \end{equation*}
        And we want to \textbf{discover some grouping} based on regularities or similarities in those data. Conceptually:
        \begin{equation*}
            \left\{x_1, x_2, \ldots, x_N\right\} \to \left\{\text{cluster}_{1}, \text{cluster}_{2}, \ldots, \text{cluster}_{K}\right\}
        \end{equation*}
        For example, suppose we have customer information but nobody tells us what kinds of customers exist. A clustering algorithm could discover groups corresponding to customers with similar behavior.

    \item \definition{Representation Learning}
        \begin{center}
            \emph{Are we learning \textbf{features/embeddings/representations} from the data?}
        \end{center}
        The primary goal here isn't necessarily to predict a category or numerical value. Instead, we want the model to \textbf{learn a useful representation of the input}. Think:
        \begin{equation*}
            x \to z
        \end{equation*}
        Where $x$ is the original input and $z$ is its learned representation. We have seen this on \autopageref{sec:unsupervised-learning-word-embedding} with word embeddings, where the model learns a vector representation of words that captures semantic meaning, or more generally in autoencoders (\autopageref{fig:autoencoder}), where the model learns a compressed representation of the input data.
\end{itemize}
These definitions tell us \textbf{what task we are trying to solve}.

\begin{table}[!htp]
    \centering
    \begin{tabular}{@{} l p{18em} @{}}
        \toprule
        \textbf{Task} & \textbf{Definition} \\
        \midrule
        \textbf{Classification} & Discrete label(s) \\
        \textbf{Regression} & Continuous value(s) \\
        \textbf{Clustering} & Discover groups, no predefined labels \\
        \textbf{Representation Learning} & Learn useful features / embeddings / representations \\
        \bottomrule
    \end{tabular}
    \caption{Summary of the four tasks: classification, regression, clustering, and representation learning.}
\end{table}

\noindent
In addition, we must also consider \textbf{learning paradigms} (\autopageref{sec:machine-learning-paradigms}), which tell us \textbf{how the learning signal/information is provided to the model during learning}. In other words, the learning paradigm tells us ``\emph{where does the learning signal come from?}''. The three main paradigms are:
\begin{itemize}
    \item \definition{Supervised Learning}: train using inputs \textbf{and desired outputs (labels)}:
        \begin{equation*}
            D = \left\{(x_1, t_1), (x_2, t_2), \ldots, (x_N, t_N)\right\}
        \end{equation*}
        The target $t_i$ tells the model what it should have predicted for input $x_i$. The model learns to map inputs to outputs.

    \item \definition{Unsupervised Learning}: train using the data \textbf{without desired outputs (no labels)}:
        \begin{equation*}
            D = \left\{x_1, x_2, \ldots, x_N\right\}
        \end{equation*}
        The model exploits regularities/structure in the data to learn useful representations or groupings. There is no explicit feedback on what the correct output should be.

    \item \definition{Reinforcement Learning}: learning happens through \textbf{interaction with an environment}:
        \begin{equation*}
            \text{state} \to \text{action} \to \text{reward/penalty}
        \end{equation*}
        The reward provides the learning signal, with the goal of maximizing long-term cumulative reward.
\end{itemize}
Now, we can analyze each statement in the exercise:
\begin{itemize}
    \item[\textcolor{Red2}{\faIcon{times}}] \textcolor{Red2}{\textbf{\emph{A model trained only to reconstruct its input is necessarily a reinforcement learning agent.}}}

        This is \textbf{false}. Reconstructing the input means something like:
        \begin{equation*}
            x \to \text{encoder} \to z \to \text{decoder} \to \hat{x}
        \end{equation*}
        With the goal of minimizing the difference between $x$ and $\hat{x}$:
        \begin{equation*}
            \hat{x} \approx x
        \end{equation*}
        That's the typical setup of an \textbf{autoencoder}.

        Instead, reinforcement learning involves an agent interacting with an environment and receiving rewards or penalties based on its actions. Reconstructing the input does not involve any interaction with an environment or reward signal. Thus, this statement is incorrect.

    \item[\textcolor{Green3}{\faIcon{check}}] \textcolor{Green3}{\textbf{\emph{A machine-health estimator may combine a discrete status label and a continuous risk score.}}}

        This one is probably among the options that made us hesitate. First of all, a \definition{Machine-Health Estimator} is simply a model that looks at data from a physical machine and estimates \textbf{how healthy the machine is}. For example, imagine an industrial motor equipped with sensors measuring:
        \begin{equation*}
            x = \left[
                \text{temperature}, \text{vibration}, \text{pressure}, \text{RPM}, \ldots
            \right]
        \end{equation*}
        A Neural Network receives those measurements and could produce \textbf{two outputs}:
        \begin{equation*}
            x \to
            \begin{cases}
                \text{status} = \text{``warning''} \\
                \text{risk score} = 0.73
            \end{cases}
        \end{equation*}
        These two outputs correspond to two different tasks:
        \begin{itemize}
            \item The \textbf{status} is a \textbf{discrete label}:
                \begin{equation*}
                    \text{status} \in \left\{\text{``healthy''}, \text{``warning''}, \text{``failure''}\right\}
                \end{equation*}
                This is a \textbf{classification task}, where the model predicts one of several discrete categories.

            \item The \textbf{risk score} is a \textbf{continuous value}:
                \begin{equation*}
                    \text{risk score} \in [0, 1]
                \end{equation*}
                This is a \textbf{regression task}, where the model predicts a continuous quantity.
        \end{itemize}
        A Neural Network can simultaneously perform both tasks, making it a multi-task learning model. Therefore, this statement is \textbf{true}.

        \textcolor{Green3}{\faIcon[regular]{lightbulb} \textbf{Tip:}} we shouldn't assume that a model can \emph{only perform one type} of task.

    \item[\textcolor{Green3}{\faIcon{check}}] \textcolor{Green3}{\textbf{\emph{Regression predicts continuous quantities such as price or temperature.}}}

        This is simply the definition of regression, so this statement is \textbf{true}.

    \item[\textcolor{Red2}{\faIcon{times}}] \textcolor{Red2}{\textbf{\emph{Clustering requires the desired class label for every training example.}}}

        This is the \textbf{opposite} of the definition of clustering, so this statement is \textbf{false}. If we have:
        \begin{equation*}
            D = \left\{(x_1, t_1), (x_2, t_2), \ldots, (x_N, t_N)\right\}
        \end{equation*}
        And $y_i$ tells us the correct class, then we are in a \textbf{supervised learning} scenario, not clustering. In clustering we instead have:
        \begin{equation*}
            D = \left\{x_1, x_2, \ldots, x_N\right\}
        \end{equation*}
        And try to discover groups such as:
        \begin{equation*}
            \left\{x_1, x_2, \ldots, x_N\right\} \to \left\{\text{cluster}_{1}, \text{cluster}_{2}, \ldots, \text{cluster}_{K}\right\}
        \end{equation*}
        No predefined class labels are required, so this statement is \textbf{false}.

    \item[\textcolor{Red2}{\faIcon{times}}] \textcolor{Red2}{\textbf{\emph{A task is unsupervised only because its input is high-dimen\break{}sional.}}}

        Classic distraction. \hl{Dimensionality has nothing to do with whether learning is supervised or unsupervised.} If we have a $1'000'000$-dimensional image and its label (e.g., ``cat'' or ``dog''), we are in a supervised learning scenario. What matters is \textbf{supervision/targets}, not dimensionality. Therefore, this statement is \textbf{false}.

    \item[\textcolor{Green3}{\faIcon{check}}] \textcolor{Green3}{\textbf{\emph{Classification predicts one or more discrete labels.}}}

        This is \textbf{true} because it is the definition of classification. If the classification predicts one label, it is called \textbf{single-label classification}. If it predicts multiple labels, it is called \textbf{multi-label classification}. In both cases, the output consists of discrete categories.

    \item[\textcolor{Green3}{\faIcon{check}}] \textcolor{Green3}{\textbf{\emph{Predicting a binary sentiment label from examples is supervised learning.}}}

        A \definition{Binary Sentiment Label} is a label describing the \textbf{sentiment of something using two possible categories}. For example, suppose we want a model to understand whether a movie review is positive or negative:
        \begin{itemize}
            \item ``\emph{This movie was fantastic!}'' $\to$ \textbf{positive}
            \item ``\emph{I didn't like this movie at all.}'' $\to$ \textbf{negative}
        \end{itemize}
        For example, suppose we want a model to understand whether a movie review is positive or negative:
        \begin{equation*}
            y \in \left\{\text{positive}, \text{negative}\right\}
        \end{equation*}
        So \textbf{binary sentiment classification} is a \textbf{classification task}: given some text, classify its sentiment into one of two categories, typically \textbf{positive} or \textbf{negative}. Since we have \textbf{labeled examples} (the text and its corresponding sentiment label), this is a \textbf{supervised learning} scenario. Therefore, this statement is \textbf{true}.

    \item[\textcolor{Green3}{\faIcon{check}}] \textcolor{Green3}{\textbf{\emph{Learning an embedding can be useful even if no semantic class label is available.}}}

        \definition{Semantic Class Label} means a label that gives an example an \textbf{explicit human-interpretable category/meaning}. For example, in a dataset of images, a semantic class label could be ``cat'', ``dog'', or ``car''. So the statement means that we might have data $x_1$, $x_2$, $\dots$, $x_{N}$, but nobody has annotated them with meaningful categories such as ``dog'', ``cat'', or ``car''. However, we can still learn an \textbf{embedding} $x \to z$, where $z$ is a useful vector representation of the input $x$. This is often done in \textbf{unsupervised learning} scenarios, where we want to learn a representation that captures important features of the data without relying on explicit labels. Therefore, this statement is \textbf{true}.
\end{itemize}